% Auto-generated by tools/qbe.py problem-latex-export.
% Problem-specific proof note for copying into a user manuscript.
% Generated: 2026-06-18 22:02:31.  Task: QBE-OP-OPTCTRL-001.
% Requires ordinary amsmath/amssymb notation.  The Lean declaration names are
% included as audit anchors and may be removed from a polished paper draft.

\subsection{A concrete block encoding of the transfer operator}

Let \(T\) be a one-qubit time register, let \(\tau\) be a one-qubit type
register, and let \(S\) be a one-qubit passive state register.  For the
concrete case \(k=1\), define
\[
  E_1
  =
  |0\rangle\langle 1|_T
  \otimes
  |0\rangle\langle 1|_\tau
  \otimes I_S .
\]
The operator \(E_1\) is a partial isometry and is not itself a unitary on the
system register.  In the ABEIS Lean development this non-unitarity obstruction
is recorded by \texttt{OptimalControl.exampleOperator\_not\_rationalOrthogonal}.

Add one block-encoding auxiliary qubit \(a\), initialized and projected in
\(|0\rangle_a\).  We use computational-basis states
\(|a,t,b,s\rangle\), where \(t,b,a\in\{0,1\}\) and
\(s\in\{0,1\}\).  Consider the logical reversible circuit
\[
  U_{\mathrm{evo}}
  =
  \bigl(X_a X_T X_\tau\bigr)
  \;\mathrm{CCX}_{\tau,T\to a},
\]
where the three \(X\) gates act on disjoint qubits and hence form one
parallel layer after the Toffoli layer.  The state register \(S\) is untouched.

\paragraph{Theorem.}
The circuit \(U_{\mathrm{evo}}\) is an exact one-ancilla block encoding of
\(E_1\):
\[
  (\langle 0|_a\otimes I_{T\tau S})
  U_{\mathrm{evo}}
  (|0\rangle_a\otimes I_{T\tau S})
  =
  E_1 .
\]
In the logical gate library \(\{X,\mathrm{CNOT},\mathrm{Toffoli}\}\),
its certified score is
\[
  (\mathrm{gateCount},\mathrm{depth},\mathrm{auxiliaryQubits},
  \mathrm{oracleCalls})
  =
  (4,2,1,0).
\]
Equivalently, it is also an \((\alpha,a,\varepsilon)=(1,1,0)\)
approximate block encoding.

\paragraph{Proof.}
We prove the claim on computational basis states; linearity then gives the
matrix identity.  The Toffoli first maps
\[
  |a,t,b,s\rangle
  \longmapsto
  |a\oplus tb,t,b,s\rangle .
\]
The parallel layer \(X_aX_TX_\tau\) then maps this state to
\[
  |1\oplus a\oplus tb,\;1\oplus t,\;1\oplus b,\;s\rangle .
\]
Thus \(U_{\mathrm{evo}}\) is unitary, because it is a composition of
reversible elementary gates.  In the finite Lean model this is the rational
orthogonality theorem
\[
  \texttt{OptimalControl.evolvedEqFlipUnitary\_isRationalOrthogonal}.
\]

Now restrict the input auxiliary to the clean value \(a=0\).  The above formula
becomes
\[
  U_{\mathrm{evo}}|0,t,b,s\rangle
  =
  |1\oplus tb,\;1\oplus t,\;1\oplus b,\;s\rangle .
\]
Projecting the output auxiliary onto \(\langle 0|_a\) keeps exactly the
branches with \(1\oplus tb=0\), namely \(t=b=1\).  The four cases are
\[
\begin{array}{c|c|c}
(t,b) & U_{\mathrm{evo}}|0,t,b,s\rangle
  & (\langle0|_a\otimes I)U_{\mathrm{evo}}|0,t,b,s\rangle \\
\hline
(0,0) & |1,1,1,s\rangle & 0 \\
(0,1) & |1,1,0,s\rangle & 0 \\
(1,0) & |1,0,1,s\rangle & 0 \\
(1,1) & |0,0,0,s\rangle & |0\rangle_T|0\rangle_\tau|s\rangle_S .
\end{array}
\]
Therefore the clean block sends \(|1\rangle_T|1\rangle_\tau|s\rangle_S\)
to \(|0\rangle_T|0\rangle_\tau|s\rangle_S\) and sends all other
\((T,\tau)\)-basis branches to zero.  This is exactly
\[
  |0\rangle\langle1|_T\otimes
  |0\rangle\langle1|_\tau\otimes I_S
  =
  E_1 .
\]
The formal Lean certificate proves the same clean-block equality entrywise:
\[
  \texttt{OptimalControl.evolvedEqFlipVerified},
  \quad
  \texttt{OptimalControl.evolvedEqFlipUnitary\_cleanBlock}.
\]
The resource count is also explicit.  The circuit has one Toffoli gate and
three \(X\) gates, hence gate count \(4\).  The Toffoli layer must precede
the flips, but the three flips act on disjoint qubits and form one parallel
layer, hence depth \(2\).  It uses one block-encoding auxiliary qubit and no
oracle calls.  This resource tuple is certified by
\[
  \texttt{OptimalControl.evolvedEqFlipCandidate\_cost}.
\]

Finally, the approximate block-encoding statement follows from exactness.
With \(\alpha=1\) and \(\varepsilon=0\), the norm error is zero.  The Lean
project records this zero-error approximate certificate as
\[
  \texttt{OptimalControl.evolvedEqFlipZeroErrorApprox}.
\]

\paragraph{Finite convergence diagnostic.}
After correcting the resource order to
\[
  \mathrm{gateCount}
  \;>\;
  \mathrm{depth}
  \;>\;
  \mathrm{auxiliaryQubits}
  \;>\;
  \mathrm{oracleCalls},
\]
the ABEIS verifier exhaustively enumerated the reduced three-bit logical gate
library containing all orientations of \(X\), CNOT, and Toffoli.  In that
finite library there is no clean-block candidate with at most three gates, no
depth-one layered candidate with at most four gates, and the depth-two witness
with four gates is exactly the construction above.  This is recorded as a
finite verifier result rather than as a Lean lower-bound theorem.

\paragraph{Scope.}
This note proves the concrete \(r=1,k=1\) logical reversible
permutation-matrix instance.  It is not a hardware-decomposed resource theorem,
not yet a general arbitrary-register construction, and not a Lean-proved
global optimality theorem.
